\begin{lemma}[Kunz]
\label{lemma-frobenius-flat-regular}
\begin{reference}
\cite{Kunz-flat}
\end{reference}
Let $p$ be a prime number.
Let $A$ be a Noetherian ring with $p = 0$.
The following are equivalent
\begin{enumerate}
\item $A$ is regular, and
\item $F : A \to A$, $a \mapsto a^p$ is flat.
\end{enumerate}
\end{lemma}

\begin{proof}
Observe that $\Spec(F) : \Spec(A) \to \Spec(A)$ is the identity map.
Being regular is defined in terms of the local rings and being flat
is something about local rings, see
Algebra, Lemma \ref{algebra-lemma-flat-localization}.
Thus we may and do assume $A$ is a Noetherian
local ring with maximal ideal $\mathfrak m$.

\medskip\noindent
Assume $A$ is regular. Let $x_1, \ldots, x_d$ be a
system of parameters for $A$. Applying $F$ we find
$F(x_1), \ldots, F(x_d) = x_1^p, \ldots, x_d^p$,
which is a system of parameters for $A$. Hence $F$ is flat, see
Algebra, Lemmas \ref{algebra-lemma-CM-over-regular-flat} and
\ref{algebra-lemma-regular-ring-CM}.

\medskip\noindent
Conversely, assume $F$ is flat. Write $\mathfrak m = (x_1, \ldots, x_r)$
with $r$ minimal. Then $x_1, \ldots, x_r$ are independent in the sense
defined above. Since $F$ is flat, we see that $x_1^p, \ldots, x_r^p$
are independent, see Lemma \ref{lemma-flat-extension-independent}.
Hence $\text{length}_A(A/(x_1^p, \ldots, x_r^p)) = p^r$ by
Lemma \ref{lemma-3}.
Let $\chi(n) = \text{length}_A(A/\mathfrak m^n)$ and recall
that this is a numerical polynomial of degree $\dim(A)$, see
Algebra, Proposition \ref{algebra-proposition-dimension}.
Choose $n \gg 0$. Observe that
$$
\mathfrak m^{pn + pr} \subset F(\mathfrak m^n)A \subset \mathfrak m^{pn}
$$
as can be seen by looking at monomials in $x_1, \ldots, x_r$. We have
$$
A/F(\mathfrak m^n)A = A/\mathfrak m^n \otimes_{A, F} A
$$
By flatness of $F$ this has length $\chi(n) \text{length}_A(A/F(\mathfrak m)A)$
(Algebra, Lemma \ref{algebra-lemma-pullback-module})
which is equal to $p^r\chi(n)$ by the above. We conclude
$$
\chi(pn + pr) \geq p^r\chi(n) \geq \chi(pn)
$$
Looking at the leading terms this implies $r = \dim(A)$, i.e., $A$ is regular.
\end{proof}